\documentclass[11pt,a4paper]{scrartcl}

\usepackage{amssymb}
\usepackage{tikz} 
\usepackage{amsmath}

\usepackage[utf8]{inputenc}
\usepackage[T1]{fontenc}
\newcommand{\ats}[1]{\par\noindent\textit{Atsakymas:} #1}
\usepackage{cancel}
\usepackage{tikz}
\usepackage{lmodern} 
\usepackage{amsmath,amssymb,amsthm}
\usepackage{graphicx}
\usepackage[dvipsnames,svgnames]{xcolor}
\usepackage{hyperref}
\usepackage{mathrsfs}
\usepackage[shortlabels]{enumitem}
\usepackage[obeyFinal,textsize=scriptsize,shadow,loadshadowlibrary]{todonotes}
\usepackage{mathtools}
\usepackage{microtype}

%%% Paragraph layout
\setlength{\parindent}{0pt}
\setlength{\parskip}{0.6\baselineskip}

%%% Colored sections
\renewcommand*{\sectionformat}%
  {\color{purple}\S\thesection\enskip}
\renewcommand*{\subsectionformat}%
  {\color{purple}\S\thesubsection\enskip}
\renewcommand*{\subsubsectionformat}%
  {\color{purple}\S\thesubsubsection\enskip}
\addtokomafont{paragraph}{\color{orange!35!black}\P\ }
\KOMAoptions{numbers=noenddot}

%%% Title
\addtokomafont{subtitle}{\Large}
\setkomafont{author}{\Large\scshape}
\setkomafont{date}{\Large\normalsize}
% Neigiamas vertikalus tarpas, kuris pakelia dokumento antraštę į viršų. 
% Galite keisti -2.5cm į -3cm (aukščiau) arba -1.5cm (žemiau).
\titlehead{\vspace*{-7.5cm}} 

% PRIDĖTA: Automatiškai perrašome datą į mokyklos pavadinimą
\AtBeginDocument{%
  \date{\normalsize Vilniaus licėjus}
}

%%% Page setup
\usepackage[headsepline]{scrlayer-scrpage}
\renewcommand{\headfont}{}
\addtolength{\textheight}{3.14cm}
\setlength{\footskip}{0.5in}
\setlength{\headsep}{10pt}
% Puslapio antraštė turi rodyti tik konkrečius dokumento duomenis.
% Nustatome juos aiškiai, kad neatsirastų "\@author" / "\@title" arba senų reikšmių.
\makeatletter
\AtBeginDocument{%
  \ihead{\footnotesize\textbf{Andrius Berniukevičius}}%
  \ohead{\footnotesize\textbf{\@title}}%
}
\makeatother
\automark{section}
\chead{}
\cfoot{\pagemark}

%%% Macros
\providecommand{\ol}{\overline}
\providecommand{\ul}{\underline}
\providecommand{\wt}{\widetilde}
\providecommand{\wh}{\widehat}
\providecommand{\eps}{\varepsilon}
\providecommand{\half}{\frac{1}{2}}
\providecommand{\inv}{^{-1}}
\newcommand{\dang}{\measuredangle} %% Directed angle
\newcommand{\vocab}[1]{\textbf{\color{blue}\sffamily #1}}
\providecommand{\CC}{\mathbb C}
\providecommand{\FF}{\mathbb F}
\providecommand{\NN}{\mathbb N}
\providecommand{\QQ}{\mathbb Q}
\providecommand{\RR}{\mathbb R}
\providecommand{\ZZ}{\mathbb Z}
\providecommand{\ts}{\textsuperscript}
\providecommand{\dg}{^\circ}
\providecommand{\ii}{\item}
\providecommand{\defeq}{\coloneqq}
\DeclareMathOperator*{\lcm}{lcm}
\DeclareMathOperator*{\argmin}{arg min}
\DeclareMathOperator*{\argmax}{arg max}
\providecommand{\hrulebar}{\par
\hspace{\fill}\rule{0.95\linewidth}{.7pt}\hspace{\fill}
\par\nointerlineskip \vspace{\baselineskip}}

%%% Optional colored theorems
\usepackage{thmtools}
\usepackage[framemethod=TikZ]{mdframed}
\mdfdefinestyle{mdbluebox}{%
  roundcorner=10pt,
  linewidth=1pt,
  skipabove=12pt,
  innerbottommargin=9pt,
  skipbelow=2pt,
  linecolor=blue,
  nobreak=true,
  backgroundcolor=TealBlue!5,
}
\declaretheoremstyle[
  headfont=\sffamily\bfseries\color{MidnightBlue},
  mdframed={style=mdbluebox},
  headpunct={\\[3pt]},
  postheadspace={0pt}
]{thmbluebox}
\mdfdefinestyle{mdredbox}{%
  linewidth=0.5pt,
  skipabove=12pt,
  frametitleaboveskip=5pt,
  frametitlebelowskip=0pt,
  skipbelow=2pt,
  frametitlefont=\bfseries,
  innertopmargin=4pt,
  innerbottommargin=8pt,
  nobreak=true,
  backgroundcolor=Salmon!5,
  linecolor=RawSienna,
}
\declaretheoremstyle[
  headfont=\bfseries\color{RawSienna},
  mdframed={style=mdredbox},
  headpunct={\\[3pt]},
  postheadspace={0pt},
]{thmredbox}
\mdfdefinestyle{mdgreenbox}{%
  skipabove=8pt,
  linewidth=2pt,
  rightline=false,
  leftline=true,
  topline=false,
  bottomline=false,
  linecolor=ForestGreen,
  backgroundcolor=ForestGreen!5,
}
\declaretheoremstyle[
  headfont=\bfseries\sffamily\color{ForestGreen!70!black},
  bodyfont=\normalfont,
  spaceabove=2pt,
  spacebelow=1pt,
  mdframed={style=mdgreenbox},
  headpunct={ --- },
]{thmgreenbox}
\mdfdefinestyle{mdblackbox}{%
  skipabove=8pt,
  linewidth=3pt,
  rightline=false,
  leftline=true,
  topline=false,
  bottomline=false,
  linecolor=black,
  backgroundcolor=RedViolet!5!gray!5,
}
\declaretheoremstyle[
  headfont=\bfseries,
  bodyfont=\normalfont\small,
  spaceabove=0pt,
  spacebelow=0pt,
  mdframed={style=mdblackbox}
]{thmblackbox}

%% Aplinkos su rėmeliais (thmtools)
\declaretheorem[style=thmbluebox,name=Teorema]{theorem}
\declaretheorem[style=thmbluebox,name=Lema]{lemma}
\declaretheorem[style=thmbluebox,name=Savybė]{property}
\declaretheorem[style=thmbluebox,name=Teiginys]{proposition}
\declaretheorem[style=thmbluebox,name=Išvada]{corollary}
\declaretheorem[style=thmbluebox,name=Teorema,numbered=no]{theorem*}
\declaretheorem[style=thmbluebox,name=Lema,numbered=no]{lemma*}
\declaretheorem[style=thmbluebox,name=Teiginys,numbered=no]{proposition*}
\declaretheorem[style=thmbluebox,name=Išvada,numbered=no]{corollary*}
\declaretheorem[style=thmgreenbox,name=Algoritmas]{algorithm}
\declaretheorem[style=thmgreenbox,name=Algoritmas,numbered=no]{algorithm*}
\declaretheorem[style=thmgreenbox,name=Teiginys]{claim}
\declaretheorem[style=thmgreenbox,name=Teiginys,numbered=no]{claim*}
\declaretheorem[style=thmredbox,name=Pavyzdys]{example}
\declaretheorem[style=thmredbox,name=Pavyzdys,numbered=no]{example*}
\declaretheorem[style=thmblackbox,name=Pastaba]{remark}
\declaretheorem[style=thmblackbox,name=Pastaba,numbered=no]{remark*}

%% Standartinės amsthm aplinkos (Apibrėžimai ir užduotys)
\theoremstyle{definition}
\ifcsname conjecture\endcsname\else\newtheorem{conjecture}{Hipotezė}\fi
\ifcsname definition\endcsname\else\newtheorem{definition}{Apibrėžimas}\fi
\ifcsname fact\endcsname\else\newtheorem{fact}{Faktas}\fi
\ifcsname answer\endcsname\else\newtheorem{answer}{Atsakymas}\fi
\ifcsname ques\endcsname\else\newtheorem{ques}{Klausimas}\fi
\ifcsname exercise\endcsname\else\newtheorem{exercise}{Užduotis}\fi
\ifcsname problem\endcsname\else\newtheorem{problem}{Uždavinys}[section]\fi
\ifcsname conjecture*\endcsname\else\newtheorem*{conjecture*}{Hipotezė}\fi
\ifcsname definition*\endcsname\else\newtheorem*{definition*}{Apibrėžimas}\fi
\ifcsname fact*\endcsname\else\newtheorem*{fact*}{Faktas}\fi
\ifcsname answer*\endcsname\else\newtheorem*{answer*}{Atsakymas}\fi
\ifcsname ques*\endcsname\else\newtheorem*{ques*}{Klausimas}\fi
\ifcsname exercise*\endcsname\else\newtheorem*{exercise*}{Užduotis}\fi
\ifcsname problem*\endcsname\else\newtheorem*{problem*}{Uždavinys}\fi

\newcounter{answerssection}
\newcommand{\answerssection}{%
  \setcounter{answerssection}{\value{section}}%
  \section{Atsakymai}%
  \setcounter{problem}{0}%
  \renewcommand{\theproblem}{\arabic{answerssection}.\arabic{problem}}%
}

%% GALUTINIS SPRENDIMAS PROOF APLINKAI
%% --- PRADŽIA: Įrodymo aplinkos sutvarkymas ---
\makeatletter

% 1. Užtikriname, kad žodis būtų "Įrodymas", net jei "babel" paketas bando jį perrašyti
\AtBeginDocument{%
  \renewcommand{\proofname}{Įrodymas}%
  \@ifpackageloaded{babel}{%
    \addto\captionslithuanian{\renewcommand{\proofname}{Įrodymas}}%
  }{}%
}

% 2. Priverstinai pašaliname scrartcl klasės bold šriftą ir paliekame TIK italic
% 3. Sujungiame su tuščiu kvadratėliu dešiniajame kampe.
\renewcommand{\qedsymbol}{\ensuremath{\Box}}
\renewenvironment{proof}[1][\proofname]{\par
  \pushQED{\qed}%
  \normalfont \topsep6\p@\@plus6\p@\relax
  \trivlist
  \item[\hskip\labelsep
        \normalfont\itshape % Priverčia naudoti tik pasvirą šriftą, kaip nuotraukoje
    #1\@addpunct{.}]\ignorespaces
}{%
  \popQED\endtrivlist\@endpefalse
}
\newenvironment{solution}{\par
  \normalfont \topsep6\p@\@plus6\p@\relax
  \trivlist
  \item[\hskip\labelsep
        \normalfont\itshape
    \textit{Sprendimas}\@addpunct{.}]\ignorespaces
}{%
  \endtrivlist\@endpefalse
}
\newcommand{\ans}[1]{\par\noindent\textit{Atsakymas:} #1}
\makeatother


\title{Geometrinis dažymas}
\author{Andrius Berniukevičius}

\newenvironment{solution}
    {\begin{list}{}{\setlength{\leftmargin}{10pt}\setlength{\rightmargin}{10pt}}\item[]}
    {\end{list}}

\begin{document}

\maketitle

\section{Nuo lyginumo prie spalvų}

Praėjusioje temoje kalbėjome apie lyginumą. Matėme, kad šachmatinis (dviejų spalvų) dažymas tėra geometrinė lyginumo išraiška: juodas langelis reiškia nelyginį skaičių, baltas – lyginį. Vienas ėjimas – viena pakeista spalva.

Šachmatinis dažymas yra tobulas įrankis, kai uždavinyje veikia $1 \times 2$ dydžio figūros (domino kauliukai) arba kai žirgas šokinėja šachmatų lentoje. Tačiau olimpiadininkai žino, kad gyvenimas neapsiriboja dalyba iš 2. 
Kas nutinka, jeigu jums reikia lentą padengti $1 \times 3$ dydžio kaladėlėmis? Arba $1 \times 4$ dydžio „laivais“? Šachmatinė lenta čia nebepadės – 2 spalvos neturi nieko bendro su figūromis, sudarytomis iš 3 ar 4 langelių.

Tokiais atvejais mes griebiamės \textbf{sudėtingesnio geometrinio dažymo}. Mes dažome erdvę 3, 4 ar net daugiau spalvų. Užuot naudoję tikrus dažus, lentelės langeliuose mes tiesiog įrašome skaičius ($0, 1, 2, 3\dots$). 

Ši technika leidžia įrodyti neįtikėtinus dalykus. Iš pradžių susipažinkime su trimis standartinėmis dažymo strategijomis, o tuomet pereisime prie tikrojo olimpiadinio kūrybiškumo.

% ==========================================
% 1 STRATEGIJA: ĮSTRIŽAINĖS
% ==========================================
\section{1 strategija: Įstrižaininis dažymas (Dalyba iš $k$)}

Kai uždavinyje figūruoja $1 \times 3$ dydžio stačiakampiai (tromino), lentą dažome 3 spalvomis (įrašome skaičius $0, 1, 2$). Tai darome \textbf{įstrižai}: kiekvienas langelis $(x, y)$ gauna spalvą, atitinkančią sumos $x+y$ liekaną padalijus iš 3.

Kokia šio dažymo magija? Ogi ta, kad kaip bepadėtume $1 \times 3$ figūrą (horizontaliai ar vertikaliai), ji \textbf{visada} uždengs lygiai po vieną kiekvienos spalvos langelį!

\begin{example}[Trūkstamas langelis]
Turime $5 \times 5$ dydžio lentą, kurioje išpjautas kairysis viršutinis kampas. Ar galima likusius $24$ langelius uždengti aštuoniomis $1 \times 3$ figūromis?
\end{example}

\begin{solution}
\textbf{Sprendimas:}
Lentos plotas yra 24, todėl teoriškai aštuonios figūros tilptų ($8 \times 3 = 24$). Tačiau pritaikykime įstrižaininį dažymą 3 spalvomis (skaičiais 0, 1, 2).

\begin{center}
\begin{tikzpicture}[scale=0.6]
    \foreach \x in {1,...,5} {
        \foreach \y in {1,...,5} {
            \pgfmathparse{int(mod(\x+\y,3))}
            \edef\c{\pgfmathresult}
            \ifnum\c=0
                \draw[fill=blue!20] (\x,-\y) rectangle ++(1,1);
                \node at (\x+0.5,-\y+0.5) {\textbf{0}};
            \else
                \draw[fill=white] (\x,-\y) rectangle ++(1,1);
                \node at (\x+0.5,-\y+0.5) {\c};
            \fi
        }
    }
    % X mark on top-left (1,1) -> coordinate in tikz is (1,-1)
    \draw[very thick, red] (1,-1) -- (2,0);
    \draw[very thick, red] (1,0) -- (2,-1);
\end{tikzpicture}
\end{center}

Pažiūrėkime į visos pilnos $5 \times 5$ lentos spalvų kiekius (suskaičiuokite skaičius brėžinyje):
\begin{itemize}
    \item Skaičius \textbf{0} pasikartoja \textbf{9} kartus.
    \item Skaičius \textbf{1} pasikartoja \textbf{8} kartus.
    \item Skaičius \textbf{2} pasikartoja \textbf{8} kartus.
\end{itemize}
Jeigu likusią lentą įmanoma uždengti $1 \times 3$ figūromis, tai kiekviena figūra uždengs lygiai po vieną $0$, $1$ ir $2$. Vadinasi, aštuonios figūros privalo uždengti lygiai aštuonius $0$, aštuonius $1$ ir aštuonius $2$. 

Kad taip nutiktų, iš lentos privalome išpjauti \textbf{0} spalvos langelį (nes nulių yra per daug). Tačiau pažiūrėkite į kairįjį viršutinį kampą – ten įrašytas skaičius \textbf{2}! Išpjovus jį, lentos spalvų balansas tampa $(9, 8, 7)$. Kadangi skaičiai nevienodi, tokios lentos padengti $1 \times 3$ figūromis neįmanoma. \hfill $\square$
\end{solution}

% ==========================================
% 2 STRATEGIJA: HIBRIDINIAI DAŽYMAI
% ==========================================
\section{2 strategija: Defektų ieškojimas ($1 \times 4$ prieš $2 \times 2$)}

Kartais visos lentos spalvų balansas yra tobulas, tačiau pati „įsibrovėlė“ figūra turi defektą.

\begin{example}[Klasikinis $2 \times 2$ pakeitimas]
Standartinė šachmatų lenta ($8 \times 8$) yra idealiai padengiama šešiolika $1 \times 4$ dydžio stačiakampių. Ar galima šią lentą padengti naudojant penkiolika $1 \times 4$ stačiakampių ir vieną $2 \times 2$ dydžio kvadratą?
\end{example}

\begin{solution}
\textbf{Sprendimas:}
Lentos plotas vis dar 64 ($15 \times 4 + 4 = 64$). Dažykime lentą įstrižai \textbf{4 spalvomis} ($0, 1, 2, 3$). Pilnoje $8 \times 8$ lentoje kiekvienos spalvos yra lygiai po $16$.

\begin{center}
\begin{tikzpicture}[scale=0.55]
    % Piešiame 5x5 fragmentą su 4 spalvų įstrižaininiu dažymu
    \foreach \x in {0,...,4} {
        \foreach \y in {0,...,4} {
            \pgfmathparse{int(mod(\x+\y, 4))}
            \edef\c{\pgfmathresult}
            \ifnum\c=0 \def\bg{white} \fi
            \ifnum\c=1 \def\bg{black!5} \fi
            \ifnum\c=2 \def\bg{black!15} \fi
            \ifnum\c=3 \def\bg{black!25} \fi
            \draw[fill=\bg, draw=gray!80] (\x,-\y) rectangle ++(1,1);
            \node at (\x+0.5,-\y+0.5) {\textbf{\c}};
        }
    }
    % Mėlynas 1x4 (tobula figūra)
    \draw[very thick, blue] (0,0) rectangle (4,-1);
    \node[blue, above, yshift = 5mm] at (2,0) {\footnotesize $1 \times 4$ surenka visas};
    
    % Raudonas 2x2 (defektas)
    \draw[very thick, red] (1,-3) rectangle (3,-1);
    \node[red, right, xshift = 10mm] at (3,-2) {\footnotesize trūksta 2};
\end{tikzpicture}
\end{center}

Kiekviena $1 \times 4$ figūra (tiek horizontali, tiek vertikali) uždengia lygiai vieną $0$, vieną $1$, vieną $2$ ir vieną $3$. Taigi, penkiolika tokių figūrų sunaudos lygiai po 15 kiekvienos spalvos langelių. Vadinasi, likęs $2 \times 2$ kvadratas \textbf{privalo} uždengti lygiai po vieną kiekvienos spalvos langelį.

Tačiau pažiūrėkime, kokias spalvas dengia bet koks $2 \times 2$ kvadratas (raudonas kontūras brėžinyje). Pagal mūsų diagonalaus dažymo formulę, matome, kad $2 \times 2$ kvadratas visada turi dvi vienodas spalvas įstrižainėje, ir \textbf{jam visada trūksta vienos spalvos} (mūsų pavyzdyje jis dengia $0, 0, 1, 3$, todėl jam trūksta $1$). 

Kadangi $2 \times 2$ kvadratas nesugeba „surinkti“ visų 4 spalvų po vieną, tokios lentos padengti neįmanoma. \hfill $\square$
\end{solution}

% ==========================================
% 3 STRATEGIJA: JUOSTINIS DAŽYMAS
% ==========================================
\section{3 strategija: Juostinis dažymas (Stripe Coloring)}

Įstrižaininis dažymas yra galingas, bet kartais figūros elgiasi ne pagal įstrižaines. Jei norime blokuoti vertikalias figūras, o horizontalioms leisti veikti laisvai, mes dažome lentą \textbf{juostomis} (eilutėmis arba stulpeliais).

\begin{example}[Ilgieji laivai]
Turime $10 \times 10$ dydžio lentą. Ar galima ją visą padengti 25 „laivais“, kurių kiekvieno dydis yra $1 \times 4$ langeliai?
\end{example}

\begin{solution}
\textbf{Sprendimas:}
Dažykime lentą 4 spalvomis, bet ne įstrižai, o \textbf{horizontaliai (juostomis)}:
1 eilutę dažome $1$, antrą $2$, trečią $3$, ketvirtą $4$, penktą vėl $1$ ir t.t.
Lentoje turime dešimt eilučių. Tai reiškia, kad spalvos 1 ir 2 pasikartos po tris kartus, o spalvos 3 ir 4 – tik po du kartus.

\begin{center}
\begin{tikzpicture}[scale=0.5]
    % Piešiame 6x6 fragmentą juostiniu dažymu
    \foreach \y in {0,...,5} {
        \pgfmathparse{int(mod(\y,4)+1)}
        \edef\c{\pgfmathresult}
        \foreach \x in {0,...,5} {
            \ifnum\c=1 \def\bg{blue!15} \fi
            \ifnum\c=2 \def\bg{green!15} \fi
            \ifnum\c=3 \def\bg{yellow!20} \fi
            \ifnum\c=4 \def\bg{red!15} \fi
            \draw[fill=\bg, draw=gray!80] (\x,-\y) rectangle ++(1,1);
            \node at (\x+0.5,-\y+0.5) {\c};
        }
    }
    % Vertikalus 1x4
    \draw[very thick, blue] (1,0) rectangle (2,-4);
    \node[blue, above, yshift=1mm] at (1.5,0) {\footnotesize V};
    
    % Horizontalus 1x4
    \draw[very thick, red] (2,-4) rectangle (6,-5);
    \node[red, right] at (6,-4.5) {\footnotesize H};
\end{tikzpicture}
\end{center}

Kadangi kiekvienoje eilutėje yra 10 langelių, iš viso $10 \times 10$ lentoje bus:
\begin{itemize}
    \item Spalvos \textbf{1}: $30$ langelių.
    \item Spalvos \textbf{2}: $30$ langelių.
    \item Spalvos \textbf{3}: $20$ langelių.
    \item Spalvos \textbf{4}: $20$ langelių.
\end{itemize}

Paanalizuokime pačius laivus. \textbf{Vertikalus (V)} laivas visada uždengs lygiai po vieną $1, 2, 3$ ir $4$. \textbf{Horizontalus (H)} laivas visas guli vienoje eilutėje, todėl jis uždengs lygiai \textbf{keturis tos pačios spalvos} langelius.

Tarkime, visoje lentoje yra $V$ vertikalių laivų. Taip pat turime horizontalių laivų: tegul $H_1$ yra skaičius horizontalių laivų, gulinčių 1-os spalvos eilutėse, o $H_3$ – gulinčių 3-ios spalvos eilutėse.
Sudarykime lygtis 1-ajai ir 3-iajai spalvai (bendras uždengtų langelių skaičius):
\[ V + 4H_1 = 30 \]
\[ V + 4H_3 = 20 \]

Atimkime antrąją lygtį iš pirmosios:
\[ 4H_1 - 4H_3 = 10 \implies H_1 - H_3 = 2{,}5 \]
Gavome visišką nesąmonę! Laivų kiekis ($H_1, H_3$) privalo būti sveikasis skaičius. Dviejų sveikųjų skaičių skirtumas niekada negali būti trupmena ($2{,}5$). Vadinasi, taip padengti lentos yra neįmanoma. \hfill $\square$
\end{solution}

% ==========================================
% AUKSINĖ TAISYKLĖ: KŪRYBIŠKUMAS
% ==========================================
\vspace{1cm}
\begin{center}
\fbox{
\begin{minipage}{0.9\textwidth}
\textbf{Auksinė taisyklė: Būkite kūrybingi – susikurkite savo dažymą!}

Šachmatinė lenta, įstrižainės ar juostos yra tik populiariausi įrankiai. Tačiau tikrasis olimpiadininkas nesusiriša su jokia taisykle – jis kuria savo dažymą pagal uždavinio poreikius! Galite dažyti kas trečią stulpelį, galite nuspalvinti tik kampus, galite lentą paversti „tinkleliu“. Svarbiausia taisyklė: jūsų dažymas yra geras, jeigu figūra (kad ir kur ją padėtumėte) visada \textbf{išlaiko tą patį spalvų invariantą} (pvz., visada dengia lyginį skaičių juodų langelių).
\end{minipage}
}
\end{center}
\vspace{0.5cm}

% ==========================================
% 4 STRATEGIJA: NESTANDARTINIS (TINKLELIS)
% ==========================================
\section{4 strategija: Nestandartinis „tinklelio“ dažymas}

\begin{example}[Maža lenta ir ilgi laivai]
Ar įmanoma $6 \times 6$ dydžio lentą uždengti devyniais $1 \times 4$ dydžio stačiakampiais?
\end{example}

\begin{solution}
\textbf{Kodėl vengiame standartinių dažymų?}
\begin{itemize}
    \item \textbf{Įstrižaininis (4 spalvų) dažymas:} Lentoje yra $36$ langeliai, todėl kiekviena iš 4 spalvų pasikartos lygiai 9 kartus. Kadangi kiekvienas $1 \times 4$ laivas visada dengia lygiai po vieną kiekvienos spalvos langelį, 9 laivai sunaudos lygiai po 9 kiekvienos spalvos langelius. Balansas tobulas, įstrižainės jokios prieštaros neranda.
    \item \textbf{Juostinis dažymas:} Jei dažytume horizontaliomis juostomis (1, 2, 3, 4, 1, 2), mes prieštarą rastume. Tačiau tam reikėtų sudaryti kintamųjų lygčių sistemą (kaip 3-ioje strategijoje) ir algebriškai įrodinėti, kad gaunasi trupmenos. 
\end{itemize}

Olimpiadose mes ieškome elegancijos. Sukurkime savo unikalų „tinklelį“, kuris problemą „nukirs“ vizualiai ir akimirksniu, be jokių lygčių! 

Nudažysime juodai tik tuos langelius, kurie \textbf{vienu metu} yra: 
2-ame, 3-iame arba 6-ame stulpelyje \textbf{IR} 2-oje, 3-ioje arba 6-oje eilutėje. 
Tokių langelių iš viso yra $3 \times 3 = \mathbf{9}$. 9 yra \textbf{nelyginis} skaičius!

\begin{center}
\begin{tikzpicture}[scale=0.6]
    % Braižome langelius
    \draw[thick] (1,-6) grid (7,0);
    
    % Spalviname sankirtas (2,3,6)
    \foreach \x in {2,3,6} {
        \foreach \y in {2,3,6} {
            \draw[fill=black!70] (\x,-\y) rectangle ++(1,1);
        }
    }
    
    % Užrašome koordinates
    \foreach \x in {1,...,6} {
        \node at (\x+0.5, 0.4) {\footnotesize \x};
        \node at (0.4, -\x+0.5) {\footnotesize \x};
    }
    
    % Parodome vieną horizontalų laivą (raudoną)
    \draw[very thick, red] (1,-2) rectangle (5,-3);
    \node[red, right] at (7,-2.5) {\footnotesize \textbf{H} (dengia 2)};
    
    % Parodome vieną vertikalų laivą (mėlyną)
    \draw[very thick, blue] (5,-1) rectangle (6,-5);
    \node[blue, right] at (7,-4.5) {\footnotesize \textbf{V} (dengia 0)};
\end{tikzpicture}
\end{center}

Ką daro bet koks padėtas $1 \times 4$ stačiakampis šioje lentoje?
\begin{itemize}
    \item \textbf{Horizontalus} laivas užims 4 gretimus stulpelius. Kad ir kur jį dėsime (pvz., per stulpelius 1-2-3-4, arba 3-4-5-6), jis \textbf{visada} kirs lygiai du mūsų „ypatinguosius“ stulpelius. Jei stačiakampis guli juodoje eilutėje, jis uždengs lygiai $\mathbf{2}$ juodus langelius. Jei baltoje – $\mathbf{0}$ juodų langelių.
    \item Pagal simetriją, \textbf{vertikalus} stačiakampis irgi visada uždengs arba $\mathbf{2}$, arba $\mathbf{0}$ juodų langelių.
\end{itemize}

\textbf{Atradimas:} Kiekvienas laivas \textit{visada} uždengia lyginį juodų langelių skaičių. Vadinasi, 9 laivai kartu gali uždengti tik \textbf{lyginį} skaičių juodų langelių. Bet lentoje juodų langelių yra $\mathbf{9}$ (nelyginis!). Lyginiais skaičiais nelyginio uždengti neįmanoma. Prieštara! \hfill $\square$

\vspace{0.3cm}
\textit{\textbf{Pastaba:} Šį uždavinį iš tikrųjų galima išspręsti ir „Kryžminiu metodu“ (kurį išmoksime sekančioje strategijoje). Jei nudažytume tik 1-ą ir 5-ą eilutes, įrodytume, kad vertikalių laivų yra lyginis skaičius. Pakartojus tai su stulpeliais, gautume, kad horizontalių laivų irgi yra lyginis skaičius (Lyginis + Lyginis = Lyginis, o mes turime 9 laivus). Mūsų sukurtas „Tinklelis“ yra genialus triukas, leidžiantis abu šiuos kryžminio įrodymo žingsnius atlikti vienu metu!}
\end{solution}

% ==========================================
% 5 STRATEGIJA: DVIGUBOS JUOSTOS IR KRYŽMINIS ĮRODYMAS
% ==========================================
\vspace{0.5cm}
\section{5 strategija: Dvigubos juostos ir kryžminis įrodymas}

Ką daryti, jeigu figūra yra ne tiesi, o „laužta“? Aukščiau naudotas „tinklelis“ čia subyrėtų, nes laužta figūra kirstų jį chaotiškai ir neužtikrintų jokio lyginumo. Tokioms figūroms matematikai naudoja stambiausią kalibrą – dvigubas juostas, atliekamas dvejomis kryptimis.

\begin{example}[Z-tetromino magija]
Ar įmanoma $10 \times 10$ dydžio lentą uždengti 25 vienetais Z formos tetromino figūrų (žr. brėžinį)?
\end{example}

\begin{solution}
\textbf{1. Šablonų mirtis} \\
Šachmatinė lenta mus apgauna – kiekviena Z figūra visada uždengia lygiai 2 juodus ir 2 baltus langelius, todėl 25 figūros idealiai padengia $50/50$ lentos balansą. Standartinės juostos taip pat nieko neranda. 

\textbf{2. Sukuriame savo įrankį: Dvigubos juostos} \\
Spalvinkime lentos stulpelius \textbf{dvigubomis juostomis}: du juodi stulpeliai, du balti, du juodi... (1-2 juodi, 3-4 balti, 5-6 juodi, 7-8 balti, 9-10 juodi). 
Visoje $10 \times 10$ lentoje dabar turime 6 juodus stulpelius (60 langelių) ir 4 baltus (40 langelių).

\vspace{0.3cm}
Panagrinėkime, kiek juodų langelių uždengia viena Z figūra tokiame rašte:
\begin{center}
\begin{tikzpicture}[scale=0.6]
    % Braižome langelius 6x5 fragmentą
    \draw[thick] (1,-5) grid (6,0);
    
    % Dvigubos juostos (1-2 juoda, 3-4 balta, 5-6 juoda)
    \foreach \x in {1,2,5,6} {
        \foreach \y in {1,...,5} {
            \draw[fill=black!30] (\x,-\y) rectangle ++(1,1);
        }
    }
    
    % Vertikali Z
    \draw[very thick, blue, fill=blue!20, opacity=0.9] (1,-1) rectangle (2,-3);
    \draw[very thick, blue, fill=blue!20, opacity=0.9] (2,-2) rectangle (3,-4);
    \draw[very thick, blue] (1,-1) -- (2,-1) -- (2,-2) -- (3,-2) -- (3,-4) -- (2,-4) -- (2,-3) -- (1,-3) -- cycle;
    \node[blue, xshift=-1.5cm,yshift = 1cm] at (2,-4.5) {\footnotesize \textbf{V} (4 juodi)};
    
    % Horizontali Z
    \draw[very thick, red, fill=red!20, opacity=0.9] (3,-1) rectangle (5,-2);
    \draw[very thick, red, fill=red!20, opacity=0.9] (4,-2) rectangle (6,-3);
    \draw[very thick, red] (3,-1) -- (5,-1) -- (5,-2) -- (6,-2) -- (6,-3) -- (4,-3) -- (4,-2) -- (3,-2) -- cycle;
    \node[red, xshift=2.2cm,yshift = 0.5cm] at (5,-3.5) {\footnotesize \textbf{H} (1 juodas)};
\end{tikzpicture}
\end{center}

\begin{itemize}
    \item \textbf{Vertikali Z (stovi 2 stulpeliuose):} Ji krenta arba ant dviejų juodų stulpelių (dengia \textbf{4} juodus langelius), arba ant dviejų baltų (dengia \textbf{0} juodų), arba ant spalvų ribos – vieno balto ir vieno juodo stulpelio (tada dengia \textbf{2} juodus). Išvada: vertikali Z figūra \textit{visada} dengia \textbf{LYGINĮ} skaičių juodų langelių.
    \item \textbf{Horizontali Z (stovi 3 stulpeliuose):} Ji kerta stulpelių spalvų ribą. Jei ji kerta ribą (Juoda, Juoda, Balta), ji uždengia \textbf{3} juodus langelius. Jei kerta ribą (Juoda, Balta, Balta) – dengia \textbf{1} juodą langelį. Išvada: horizontali Z figūra \textit{visada} dengia \textbf{NELYGINĮ} skaičių juodų langelių!
\end{itemize}

Kadangi bendras juodų langelių skaičius lentoje yra \textbf{60 (lyginis)}, o vertikalios figūros įneša tik lyginius priedus, tai reiškia, kad bendrą lyginumą lemia horizontalios figūros. Kad jų sukurta nelyginių skaičių suma taptų lyginė (60), pačių \textbf{horizontalių Z figūrų privalo būti lyginis skaičius!} (Tarkime, $H$ yra lyginis).

\vspace{0.5cm}
\textbf{3. Kryžminis įrodymas} \\
Dabar pakeiskime dažymą: pasukime lentą 90 laipsnių kampu. Dažykime \textbf{eilutes} dvigubomis juostomis (1-2 eilutė juoda, 3-4 balta ir t.t.). Viskas apsiverčia veidrodiškai!
Dabar juodų langelių skaičius vėl lygus 60 (lyginis). Tačiau dabar horizontali Z dengia lyginį skaičių juodų langelių, o vertikali – nelyginį. Vadinasi, kad suma būtų lyginė, \textbf{vertikalių Z figūrų privalo būti lyginis skaičius!} ($V$ yra lyginis).

\textbf{4. Išvada} \\
Įrodėme, kad horizontalių figūrų skaičius ($H$) yra lyginis, vertikalių ($V$) – taip pat lyginis. 
Vadinasi, bendras figūrų skaičius privalo būti jų suma:
\[ H + V = Lyginis + Lyginis = \mathbf{LYGINIS \ SKAIČIUS}. \]
Tačiau mes turime lygiai \textbf{25} figūras, o 25 yra \textbf{NELYGINIS} skaičius! 
Gauname prieštarą. \hfill $\square$
\end{solution}

% ==========================================
% UŽDAVINIAI
% ==========================================
\newpage
\section{Uždaviniai savarankiškam sprendimui (30 iššūkių)}

Sprendžiant šiuos uždavinius naudokite geometrinį dažymą: spalvinkite langelius ir skaičiuokite figūrų padengiamas spalvas.

\begin{enumerate}
    \renewcommand{\labelenumi}{\textbf{\theenumi.}}
    \item Iš standartinės šachmatų lentos ($8 \times 8$) išpjovėme du langelius: patį apatinį kairįjį ir patį viršutinį dešinįjį. Ar galima likusius 62 langelius uždengti $1 \times 2$ dydžio domino kauliukais?
    \item Lentoje $5 \times 5$ išpjautas pats centrinis langelis. Ar įmanoma likusius 24 langelius uždengti domino kauliukais ($1 \times 2$)?
    \item Ar galima stačiakampę lentą $7 \times 7$ uždengti $1 \times 3$ dydžio tromino figūromis?
    \item Turime $7 \times 7$ dydžio lentą. Į patį centrą padedame vieną mažą $1 \times 1$ dydžio langelį. Ar galima likusią lentos dalį padengti 16 tromino figūrų ($1 \times 3$)?
    \item Ar galima $6 \times 6$ dydžio lentą padengti vienuolika $1 \times 3$ stačiakampių ir vienu „kampeliu“ (L formos iš 3 langelių)?
    \item Standartinėje $8 \times 8$ šachmatų lentoje stovi 15 vnt. T formos tetromino (4 langeliai) ir vienas $2 \times 2$ dydžio kvadratas. Įrodykite, kad taip būti negali.
    \item Kubą $3 \times 3 \times 3$ sudaro 27 maži kubeliai. Pelė pragriaužė ir suvalgė patį centrinį kubelį. Ar gali ji suvalgyti visus likusius 26 kubelius, kiekvienu žingsniu pereidama tik į gretimą kubelį (turintį bendrą sieną)?
    \item Ar galima padengti $12 \times 12$ dydžio lentą $1 \times 4$ dydžio stačiakampiais?
    \item Turime $9 \times 9$ lentą. Ją padengėme 27 figūromis $1 \times 3$. Įrodykite, kad tarp jų negali būti lygiai 14 horizontalių ir 13 vertikalių figūrų.
    \item Ar galima $6 \times 6$ lentą visiškai padengti 12 L-tromino figūrų (3 langelių kampeliais)?
    \item Ar galima $14 \times 14$ lentą visiškai padengti 49 Z-tetromino figūromis?
    \item Iš $8 \times 8$ lentos išpjovėme vieną langelį. Kokiose koordinatėse jis privalo būti, kad likusią lentos dalį būtų teoriškai įmanoma padengti 21 figūra $1 \times 3$?
    \item Iš $10 \times 10$ lentos išpjautas centrinis $2 \times 2$ kvadratas. Ar galima likusią dalį (96 langelius) padengti 24 T formos tetromino figūromis?
    \item Muziejaus salė turi formą stačiakampio, kurio matmenys $11 \times 12$. Grindis dengia plytelės $1 \times 4$. Viena plytelė sudužo, bet meistras atnešė naują plytelę, kurios dydis $2 \times 2$. Ar galima perdėlioti senas plyteles ir įkomponuoti naująją, kad grindys vėl būtų pilnai uždengtos?
    \item Plokštumoje nubrėžtas daugiakampis. Visi jo kampai yra statūs, o kiekvienos kraštinės ilgis yra nelyginis skaičius (1, 3, 5...). Įrodykite, kad šio daugiakampio ploto negalima padengti domino kauliukais $1 \times 2$.
    \item Ar galima padengti $10 \times 10$ lentą dvidešimt penkiomis L formos tetromino (4 langelių) figūromis?
    \item Turime $8 \times 8$ lentą. Prie jos pridėtas papildomas $1 \times 1$ langelis prie vieno iš kampų, taip kad bendras plotas taptų 65. Ar galima šią naują figūrą padengti trylika $1 \times 5$ stačiakampių?
    \item Ar galima erdvėje esantį bloką $6 \times 6 \times 6$ (susidedantį iš 216 kubelių) pilnai padengti plytomis, kurių matmenys yra $1 \times 2 \times 4$?
    \item $7 \times 7$ lentoje padėta 16 T-tetromino figūrų, kurios niekur nepersidengia, ir liko vienas tuščias langelis. Įrodykite, kad tas tuščias langelis privalo būti pačiame lentos centre.
    \item Įrodykite, kad jeigu lentą $m \times n$ galima padengti stačiakampiais $1 \times a$, tai bent vienas iš skaičių $m$ arba $n$ privalo dalytis iš $a$.
    \item Iš šachmatų lentos $8 \times 8$ iškirptas vienas langelis, ir lenta liko susijusi (nesubyrėjo į dalis). Ar gali būti taip, kad likusią lentos dalį galima padengti L formos tromino (iš 3 langelių) figūromis?
    \item Ar galima $9 \times 9$ lentos langelius nuspalvinti trimis spalvomis taip, kad kiekviename $2 \times 2$ kvadrate visos 3 spalvos būtų panaudotos?
    \item Iš $10 \times 10$ lentos iškirptas centrinis $4 \times 4$ kvadratas. Ar galima likusią figūrą (84 langelius) padengti $1 \times 4$ stačiakampiais?
    \item $13 \times 13$ lenta uždengta $2 \times 2$ ir $3 \times 3$ kvadratais (be persidengimų, padengta pilnai). Įrodykite, kad $2 \times 2$ kvadratų yra bent vienas.
    \item Turime kryžiaus formos penkmino figūrą (5 langeliai pliuso forma). Ar galima 20 tokių kryžių pilnai ir be persidengimų padengti $10 \times 10$ lentą?
    \item Erdvė $10 \times 10 \times 10$ suformuota iš mažų kubelių $1 \times 1 \times 1$. Ją visą norima padengti kaladėlėmis $1 \times 1 \times 4$. Įrodykite, kad tai padaryti neįmanoma.
    \item Duota $23 \times 23$ dydžio kvadrato formos lenta. Ji padengta lygiai vienu $1 \times 1$ kvadratėliu ir daugybe $2 \times 2$ bei $3 \times 3$ kvadratų. Įrodykite, kad $1 \times 1$ kvadratėlis negali būti lentos centre.
    \item Lenta $10 \times 10$ uždengta dvidešimt penkiais $1 \times 4$ stačiakampiais. Jūs išėmėte vieną horizontalų ir vieną vertikalų stačiakampį. Jums davė du naujus vertikalius stačiakampius $1 \times 4$. Ar įmanoma jais užpildyti atsiradusias skyles (galimai perdėliojant kitus likusius stačiakampius)?
    \item Lentoje surašyti nuliai. Vienu ėjimu leidžiama pasirinkti bet kokį $1 \times 3$ (horizontalią ar vertikalią) sritį ir visus joje esančius tris skaičius padidinti 1. Ar po kelių ėjimų įmanoma gauti lentą, kurios centrinis langelis yra $2026$, o visi likę langeliai aplink jį yra lygūs $2025$?
    \item Stačiakampis sudarytas iš mažesnių stačiakampių. Žinoma, kad kiekvieno iš tų mažesnių stačiakampių bent viena kraštinė (plotis arba ilgis) yra sveikasis skaičius (pvz., $3 \times 1{,}5$ arba $0{,}7 \times 4$). Įrodykite, kad viso didžiojo stačiakampio bent viena kraštinė taip pat privalo būti sveikasis skaičius.
\end{enumerate}

\end{document}